\documentclass[12pt, a4paper, oneside]{ctexart}
\usepackage{amsmath, amsthm, amssymb, bm, color, framed, graphicx, hyperref, mathrsfs}

\title{\textbf{4月17日作业}}
\author{韩岳成 524531910029}
\date{\today}
\linespread{1.5}
\definecolor{shadecolor}{RGB}{241, 241, 255}
\newcounter{problemname}
\newenvironment{problem}{\begin{shaded}\stepcounter{problemname}\par\noindent\textbf{题目\arabic{problemname}. }}{\end{shaded}\par}
\newenvironment{solution}{\par\noindent\textbf{解答. }}{\par}
\newenvironment{note}{\par\noindent\textbf{题目\arabic{problemname}的注记. }}{\par}
\everymath{\displaystyle}

\begin{document}

\maketitle

\begin{problem}
自由粒子能量为 $E$，在台阶势垒
\[ V(x) = \begin{cases} V_0 > 0, & x > 0 \\ 0, & x \leqslant 0 \end{cases} \]
$x = 0$ 处反射，粒子从左边入射，计算当粒子能量 $E > V_0$ 和 $0 < E < V_0$ 这两种情况下的反射系数和透射系数。
\end{problem}

\begin{solution}
情况一：当 $E > V_0$ 时：

粒子在区域 I ($x \leqslant 0$) 的波函数为 $\psi_1(x) = e^{ik_1 x} + r e^{-ik_1 x}$，其中 $k_1 = \sqrt{2mE}/\hbar$。

粒子在区域 II ($x > 0$) 的波函数为 $\psi_2(x) = t e^{ik_2 x}$，其中 $k_2 = \sqrt{2m(E-V_0)}/\hbar$。
由 $x=0$ 处波函数和一阶导数连续：

\[\psi_1(0) = \psi_2(0) \implies 1 + r = t\]
\[\psi_1'(0) = \psi_2'(0) \implies ik_1(1-r) = ik_2 t\]

联立解得：\[r = \frac{k_1 - k_2}{k_1 + k_2},\quad t = \frac{2k_1}{k_1 + k_2}\]

透射系数和反射系数由概率流密度比值定义：

反射系数 $R = \frac{|j_{\text{ref}}|}{|j_{\text{in}}|} = |r|^2 = \left(\frac{k_1 - k_2}{k_1 + k_2}\right)^2 = \frac{(\sqrt{E} - \sqrt{E - V_0})^2}{(\sqrt{E} + \sqrt{E - V_0})^2}$。

透射系数 $T = \frac{j_{\text{trans}}}{|j_{\text{in}}|} = \frac{k_2}{k_1} |t|^2 = \frac{4k_1 k_2}{(k_1 + k_2)^2} = \frac{4\sqrt{E(E - V_0)}}{(\sqrt{E} + \sqrt{E - V_0})^2}$。

由于 $R + T = 1$，结果自洽。

情况二：当 $0 < E < V_0$ 时：

此时 $k_2$ 变为纯虚数，即 $k_2 = i\kappa$，其中 $\kappa = \sqrt{2m(V_0-E)}/\hbar$。

此时区域 II 中的波函数为指数衰减形式 $\psi_2(x) = t e^{-\kappa x}$，概率流守恒表明透射概率流密度为0，透射系数 $T = 0$。

反射振幅 $r = \frac{k_1 - i\kappa}{k_1 + i\kappa}$，$R = |r|^2 = \frac{k_1^2 + \kappa^2}{k_1^2 + \kappa^2} = 1$ 。
\end{solution}

\begin{problem}
质量为 $m$ 的粒子在一维势场 \[V(x) = \begin{cases} 0, & 0 \leqslant x \leqslant a \\ \infty, & x > a \text{ 或 } x < 0 \end{cases}\] 中运动，粒子能量本征值 \[E_n = n^2\pi^2\hbar^2/2ma^2 \, (n = 1, 2, \cdots)\]能量本征函数 \[\psi_n(x) = \begin{cases} \sqrt{\frac{2}{a}} \sin \frac{n\pi x}{a}, & 0 \leqslant x \leqslant a \\ 0, & x > a \text{ 或 } x < 0 \end{cases}\]设 $t=0$ 时粒子归一化波函数为 \[\psi(x, 0) = \begin{cases} Ax(a-x), & 0 \leqslant x \leqslant a \\ 0, & x > a \text{ 或 } x < 0 \end{cases}\]其中常数 $A, a$ 已知，求：

(1) 粒子在 $\psi(x, 0)$ 态上能量平均值 $\bar{E}$；

(2) 粒子在 $\psi(x, 0)$ 态上的坐标与动量的平均值 $\bar{x}、\bar{p}_x$；

(3) $t=0$ 时刻动量的方均根偏差 $\Delta p_x = \sqrt{\langle \hat{p}_x^2 \rangle - \langle \hat{p}_x \rangle^2}$ ($\langle \hat{p}_x \rangle \equiv \bar{p}_x$)；

(4) 任意 $t > 0$ 时刻粒子波函数 $\psi(x, t)$ 表达式。
\end{problem}

\begin{solution}
因为 $t=0$ 时 $\psi(x, 0)$ 是归一化波函数，因此 \[\int_0^a A^2 x^2 (a-x)^2 dx = 1 \implies A^2 \frac{a^5}{30} = 1 \implies A = \sqrt{\frac{30}{a^5}}\]

(1) 能量平均值 $\bar{E} = \langle \hat{H} \rangle = \int_0^a \psi(x, 0)^+ \hat{H} \psi(x, 0) dx$.\[\hat{H} \psi = -\frac{\hbar^2}{2m} \frac{d^2}{dx^2} (A(ax-x^2)) = \frac{\hbar^2}{m} A\]
\[ \Rightarrow \bar{E} = \frac{\hbar^2 A}{m} \int_0^a A(ax-x^2) dx = \frac{\hbar^2 A^2}{m} \frac{a^3}{6} = \frac{5\hbar^2}{m a^2}\]

(2) 坐标平均值 $\bar{x} = \int_0^a x |\psi(x,0)|^2 dx = \int_0^a A^2 x^3(a-x)^2 dx = A^2 (\frac{a^6}{4} - \frac{2a^6}{5} + \frac{a^6}{6}) = A^2 \frac{a^6}{60} = \frac{a}{2}$。
动量平均值 $\bar{p}_x = \int_0^a \psi(x,0)^+ (-i\hbar \frac{d}{dx}) \psi(x,0) dx = -i\hbar A^2 \int_0^a x(a-x)(a-2x) dx = 0$。

(3) 动量的方均根偏差 $\Delta p_x$：
由于 $\langle \hat{p}_x^2 \rangle = 2m\langle \hat{H} \rangle = \frac{10\hbar^2}{a^2}$（因为在方势阱中，动能即总能量，势能为0），且 $\bar{p}_x = 0$，所以 $\Delta p_x = \sqrt{\langle \hat{p}_x^2 \rangle} = \frac{\sqrt{10}\hbar}{a}$。

(4) 将初始态按能量本征态展开：$\psi(x, 0) = \sum_{n=1}^\infty c_n \psi_n(x)$。

展开系数 $c_n = \int_0^a \psi_n^*(x)\psi(x,0)dx = \sqrt{\frac{2}{a}} A \int_0^a x(a-x) \sin \frac{n\pi x}{a} dx$。
计算积分，运用分部积分法得到 $c_n = \frac{8\sqrt{15}}{n^3 \pi^3} [1 - (-1)^n] / 2$。
当 $n$ 为偶数时，$c_n = 0$；
当 $n$ 为奇数时，$c_n = \frac{8\sqrt{15}}{n^3 \pi^3}$。
所以 $t$ 时刻的波函数为 $\psi(x, t) = \sum_{n=1,3,5\dots}^\infty \frac{8\sqrt{15}}{n^3 \pi^3} \sqrt{\frac{2}{a}} \sin\left(\frac{n\pi x}{a}\right) e^{-iE_n t/\hbar}$。
\end{solution}

\begin{problem}
计算算符对易关系：$[x, L_x], [p_x, L_y], [L_x, L_y], [L^2, L_x]$；计算氢原子的哈密顿量算符与 $L^2$ 算符之间的对易关系。
\end{problem}

\begin{solution}
(1) $[x, L_x] = [x, yp_z - zp_y] = y[x, p_z] - z[x, p_y] = 0$

(2) $[p_x, L_y] = [p_x, zp_x - xp_z] = z[p_x, p_x] + [p_x, z]p_x - x[p_x, p_z] - [p_x, x]p_z = -(-i\hbar)p_z = i\hbar p_z$

(3) $[L_x, L_y] = [yp_z - zp_y, zp_x - xp_z] = y[p_z, zp_x] - x[yp_z, p_z] - [zp_y, zp_x] + [zp_y, xp_z]$
其中含有不为零的对易子的项为：$y[p_z, z]p_x + x[z, p_z]p_y = y(-i\hbar)p_x + x(i\hbar)p_y = i\hbar (xp_y - yp_x) = i\hbar L_z$

(4) $[L^2, L_x] = [L_x^2 + L_y^2 + L_z^2, L_x] = [L_y^2, L_x] + [L_z^2, L_x]$
$= L_y[L_y, L_x] + [L_y, L_x]L_y + L_z[L_z, L_x] + [L_z, L_x]L_z$
$= L_y(-i\hbar L_z) + (-i\hbar L_z)L_y + L_z(i\hbar L_y) + (i\hbar L_y)L_z = 0$

(5) 氢原子的哈密顿量为 $\hat{H} = \frac{p^2}{2\mu} + V(r)$。
由于 $[L^2, r^2] = 0$（及对任何只依赖于径向 $r$ 的函数），则 $[L^2, V(r)] = 0$。
由于 $[L^2, p^2] = 0$，则 $[L^2, \hat{H}] = [\hat{H}, L^2] = 0$。
\end{solution}

\begin{problem}
设粒子的波函数 $\Psi(\phi) = \sqrt{\frac{1}{3\pi}} + \sqrt{\frac{1}{6\pi}} e^{-i\phi}$，证明该波函数归一化。对处于该状态的粒子进行 $\hat{L}_z$ 测量，可能得到哪些值？得到每个值的概率是多少？计算粒子在这个态的 $\langle \hat{L}_z \rangle$ 数值。
\end{problem}

\begin{solution}
(1) 验证归一化条件：
$\int_0^{2\pi} |\Psi(\phi)|^2 d\phi = \int_0^{2\pi} \left( \frac{1}{3\pi} + \frac{1}{6\pi} + \frac{\sqrt{2}}{6\pi} (e^{i\phi} + e^{-i\phi}) \right) d\phi$
$= \left(\frac{1}{3\pi} + \frac{1}{6\pi}\right) \times 2\pi + 0 = \frac{1}{2\pi} \times 2\pi = 1$。
故该波函数归一化。

(2) 进行 $\hat{L}_z$ 测量的可能值：
将波函数展开为 $\hat{L}_z$ 在 $\phi$ 表象中的本征态 $\Phi_m(\phi) = \frac{1}{\sqrt{2\pi}} e^{im\phi}$（其本征值为 $L_z = m\hbar$）。
$\Psi(\phi) = \frac{\sqrt{2}}{\sqrt{3}}\frac{1}{\sqrt{2\pi}} + \frac{1}{\sqrt{3}}\frac{1}{\sqrt{2\pi}} e^{-i\phi} = \sqrt{\frac{2}{3}}\Phi_0(\phi) + \sqrt{\frac{1}{3}}\Phi_{-1}(\phi)$。
因此，测量 $\hat{L}_z$ 可能得到的值为对应的本征值：$m=0$ 时的 $0$，和 $m=-1$ 时的 $-\hbar$。

(3) 对应概率分别为绝对值的平方：
得到 $0$ 的概率 $P(0) = \left|\sqrt{\frac{2}{3}}\right|^2 = \frac{2}{3}$。
得到 $-\hbar$ 的概率 $P(-\hbar) = \left|\sqrt{\frac{1}{3}}\right|^2 = \frac{1}{3}$。

(4) 期望值 $\langle \hat{L}_z \rangle$：
$\langle \hat{L}_z \rangle = \sum P_m L_{z_m} = \frac{2}{3} \times 0 + \frac{1}{3} \times (-\hbar) = -\frac{1}{3}\hbar$。
也可以通过积分直接计算 $\int \Psi^* \left(-i\hbar \frac{\partial}{\partial \phi}\right) \Psi d\phi = -\frac{1}{3}\hbar$。
\end{solution}

\begin{problem}
频率为 $\omega$ 的一维谐振子 $t=0$ 时波函数为
\[ \Psi(x, 0) = \sqrt{\frac{1}{2}} \psi_0(x) + i\sqrt{\frac{1}{5}} \psi_1(x) + A\psi_2(x) \]
($i$ 为虚数单位)，$A$ 为待定常数，$\psi_n(x)$ 为谐振子能量本征函数，求：

(1) $t$ 时刻的归一化波函数；

(2) $t$ 时刻测量能量为 $3\hbar\omega/2$ 的概率；

(3) $t$ 时刻体系的平均能量。
\end{problem}

\begin{solution}
(1) $t=0$ 时的波函数归一化：
$\int |\Psi(x,0)|^2 dx = \frac{1}{2} + \frac{1}{5} + |A|^2 = 1$。
则 $|A|^2 = 1 - 0.7 = 0.3 = \frac{3}{10}$，所以可取 $A = \sqrt{\frac{3}{10}}$。
因为谐振子的能量本征值为 $E_n = \left(n + \frac{1}{2}\right)\hbar\omega$，根据随时间演化的因子 $e^{-iE_nt/\hbar}$：
$t$ 时刻的归一化波函数为
$\Psi(x, t) = \sqrt{\frac{1}{2}} \psi_0(x) e^{-i\omega t/2} + i\sqrt{\frac{1}{5}} \psi_1(x) e^{-i3\omega t/2} + \sqrt{\frac{3}{10}} \psi_2(x) e^{-i5\omega t/2}$。

(2) $3\hbar\omega/2$ 是 $n=1$ 的能量本征态，因为能量是定态随时间演化的守恒量。
所以 $t$ 时刻测量出 $E = 3\hbar\omega/2$ 的概率等于它在 $t=0$ 时的投影模方：
$P(E_1) = |\langle \psi_1 | \Psi(x,t) \rangle|^2 = \left|i\sqrt{\frac{1}{5}} e^{-i3\omega t/2}\right|^2 = \frac{1}{5}$。

(3) 系统的平均能量（也是守恒量，与时刻 $t$ 无关）：
由于哈密顿量 $\hat{H}$ 不外显含时，
$\langle \hat{H} \rangle = \sum_n P(E_n) E_n = P(E_0) E_0 + P(E_1) E_1 + P(E_2) E_2$
$= \frac{1}{2} \left(\frac{1}{2}\hbar\omega\right) + \frac{1}{5} \left(\frac{3}{2}\hbar\omega\right) + \frac{3}{10} \left(\frac{5}{2}\hbar\omega\right)$
$= \frac{1}{4}\hbar\omega + \frac{3}{10}\hbar\omega + \frac{3}{4}\hbar\omega = \frac{10+12+30}{40}\hbar\omega = \frac{52}{40}\hbar\omega = \frac{13}{10}\hbar\omega = 1.3\hbar\omega$。
\end{solution}
\end{document}